\documentclass[11pt,a4paper]{ltjsarticle}
\usepackage[margin=24mm]{geometry}
\usepackage[haranoaji]{luatexja-preset}
\usepackage{amsmath,amssymb,amsthm,mathtools}
\usepackage{booktabs,longtable,array}
\usepackage{xcolor}
\usepackage{hyperref}
\hypersetup{colorlinks=true,linkcolor=blue!45!black,urlcolor=blue!45!black,
  pdftitle={射影四空間のPoisson構造の一次ジェット還元},
  pdfauthor={Poisson P4 research project}}
\setlength{\emergencystretch}{3em}
\newcommand{\CC}{\mathbf C}
\newcommand{\QQ}{\mathbf Q}
\newcommand{\PP}{\mathbf P}
\renewcommand{\AA}{\mathbf A}
\newcommand{\GG}{\mathbf G}
\newcommand{\cP}{\mathcal P_4}
\newcommand{\dd}{\mathrm d}
\DeclareMathOperator{\Sp}{Sp}
\DeclareMathOperator{\CSp}{CSp}
\DeclareMathOperator{\GL}{GL}
\DeclareMathOperator{\PGL}{PGL}
\DeclareMathOperator{\rank}{rank}
\DeclareMathOperator{\tr}{tr}
\DeclareMathOperator{\Stab}{Stab}
\DeclareMathOperator{\Spec}{Spec}
\DeclareMathOperator{\grad}{grad}
\newtheorem{theorem}{定理}[section]
\newtheorem{proposition}[theorem]{命題}
\newtheorem{lemma}[theorem]{補題}
\newtheorem{corollary}[theorem]{系}
\theoremstyle{remark}
\newtheorem{remark}[theorem]{注意}
\renewcommand{\proofname}{証明}
\title{射影四空間の Poisson 構造の一次ジェット還元\\
\large Toric 射影の次数三、五つの一般延長と無限遠の障害}
\author{Poisson P4 研究プロジェクト・内部研究草稿}
\date{2026年10月31日・v28-s61}
\begin{document}
\maketitle
\begin{abstract}
射影四空間の一般階数四 Poisson 構造を捉える既存の61変数模型に対し、
primitive な一次ジェットの16次元空間への射影を調べる。
既知の toric 成分 $T01$ と楕円型成分 $T39=q_{5,1}$、$T40=q_{5,2}$ は、
いずれもこの射影で優越的かつ一般に有限である。
全134式の Jacobian を用いて各成分に滑らかな非分岐点を与え、十分一般の同じ一次項には
少なくとも五つの相異なる、孤立した被約な延長が存在することを示す。
さらに、五超平面が射影標構をなす toric 開部分の全延長を、
16変数の明示方程式の解の $S_5$ 軌道として表し、さらにラベルを除いた20変数・96二次式へ還元する。
普遍的一次消去と有限な近傍の構成により、toric成分の一般次数が3、順序付きcoframeの次数が360であることを証明する。
全S61の未飽和斉次化には、非零一次項ごとに3次元の旗の無限遠部分が必ず現れることも示す。
基礎となる perspective 正規化は有限の多項式公式で与えられ、一意である。
特殊層では、係数階数一の5次元軌道上の全解を45変数・134二次式の固定 fibre へ帰着し、
二つの既知28次元成分の共通一次像を閉既約9次元集合として決定する。
MPSの全40図式を一次資料に照合すると、既知51個の一般階数四成分のうち
一般一次項上に残る成分は上述の三つに限られる。
未知成分、全体の一般 fibre の尽くし、特殊 fibre は未解決であり、成分数の下界57は変更しない。
本稿は公表上の新規性の調査を完了していない内部研究草稿である。
\end{abstract}

\section{対象、既存模型、主張の範囲}
基礎体を $\CC$ とし、対象を
\[
 \cP=\{[\pi]\in\PP H^0(\PP^4,\wedge^2T_{\PP^4}):[\pi,\pi]=0\}
\]
とする。これは非零スカラー倍を除いた Jacobi スキームであり、$\PGL_5$ 商は取らない。
研究全体の既知の成分数の下界は57であるが、上界は未確定である。

既存の第2版\cite{v2}は、一般階数四の全既約成分に対応する正の重み付き模型
\[
 S_{61}=\Spec\CC[u_1,\ldots,u_{61}]/I_{61}
\]
を構成する。重みは $1^{16},2^{45}$、保存された定義式は重み4の70式と重み5の64式である。
一次項を零にする45変数断面の分類を、$S_{61}$ 全体の分類と同一視しない。
本稿はこの模型の重み1部分への射影と、その一般 fibre および特定の特殊層を扱う。
以下で成分の射影を論じるときはまず被約成分を対象とし、全134式のスキーム構造を
用いる箇所は明記する。

\subsection{成分対応に必要な幾何}
固定点 $p=[1:0:0:0:0]$ で評価が非退化な開部分を $\mathcal U\subset\cP$ とする。
一般階数四の各既約成分は $\mathcal U$ と交わる。
実際、連結群 $\PGL_5$ は既約成分を保ち、任意の非退化点を $p$ に移せる。
従って $\mathcal U$ の既約成分は、全空間の一般階数四成分と一対一に対応する。

評価の射影値の空間は $\GL_4/\CSp_4$、次元は5である。
その安定化群 $\CSp_4$ は連結なので、associated bundle の fibre の成分を置換しない。
定数項を
\[
 J=\partial_1\wedge\partial_2+\partial_3\wedge\partial_4
\]
に固定した fibre に移り、4次元の perspective 群を除く。
第2版の55変数の三角消去と20個の Bianchi 恒等式による ideal の同一性を用いると、
対応する成分 $C,C_{61}$ の次元は
\[
 \dim C=5+4+\dim C_{61}=9+\dim C_{61}
\]
を満たす。三角消去の分母は非零の有理定数のみであり、形式近傍に限定した主張ではない。
この成分対応の使用は、保存された exact elimination の成立を入力とする。
本稿の一次ジェットの議論だけで、134式の全構築や全既約分解を置換するものではない。

\section{Perspective 作用の有限公式}
$U=\CC^4$、affine 座標 $y=(y_1,\ldots,y_4)^t$、Euler 場 $E=y$ とする。
以下では同じ記号を反対称係数行列にも用い、$v\wedge w=vw^t-wv^t$ と約束する。
特に
\[
 J=\begin{pmatrix}0&1&0&0\\-1&0&0&0\\0&0&0&1\\0&0&-1&0\end{pmatrix}.
\]
射影空間上の bivector の $x_0=1$ chart への制限を
\begin{equation}\label{eq:decomposition}
 P=J+L+Q+C
\end{equation}
と書く。係数次数はそれぞれ0,1,2,3であり、射影延長条件により
$C=E\wedge B$、$B$ は二次係数ベクトル場である。
これは、斉次二次 bivector を affine 座標に射影した際、三次項が
$\partial_0$ を含む部分から生じることによる。

$a\in U^*$ を係数列で表し、射影変換
\[
 x'_0=x_0+a^t x,\qquad x'_i=x_i\ (1\le i\le4)
\]
から生じる $g_a(y)=y/(1+a^t y)$ を考える。
以下の公式は actual pushforward に対するものであり、Lie derivative の規約とは
無限小で符号が逆になる。

\begin{proposition}[有限の perspective 公式]\label{prop:perspective}
$f=a^t y$、$h=1-f$ と置くと、\eqref{eq:decomposition} に対して
\begin{align}
 (g_a)_*P={}&h^2(J+E\wedge Ja)
       +h(L+E\wedge La)+Q+E\wedge Qa+C.\label{eq:finite}
\end{align}
特に係数次数ごとの変換は
\begin{align}
 L'&=L-2fJ+E\wedge Ja,\label{eq:Ltransform}\\
 Q'&=Q-fL+E\wedge La+f^2J-2fE\wedge Ja,\label{eq:Qtransform}\\
 C'&=C+E\wedge Qa-fE\wedge La+f^2E\wedge Ja.\label{eq:Ctransform}
\end{align}
すべての係数は元の係数と $a$ の多項式である。
\end{proposition}
\begin{proof}
$g_a^{-1}(y)=y/h$ における微分は $h(I-ya^t)$ である。
任意の反対称行列 $M$ について
\[
 (I-ya^t)M(I-ay^t)=M+E\wedge Ma
\]
が成立する。$M=C=E\wedge B$ なら右辺は $hC$ である。
$P(y/h)=J+h^{-1}L+h^{-2}Q+h^{-3}C$ に左右の微分を掛けると
三次項の分母も相殺し、\eqref{eq:finite}を得る。次数別に整理すれば残りも従う。
\end{proof}

$g_ag_b=g_{a+b}$、$g_a^{-1}=g_{-a}$ である。
$a$ の重みを1とすれば、$L,Q,C$ の重み1,2,3を保つ多項式作用である。
公式は affine chart の有理式から導いたが、両辺が射影大域切断に由来するので
chart の稠密開集合での一致から大域的にも一致する。

\section{Primitive gauge と16次元の残存表現}
一次 bivector に対して、一次形式
\[
 \tau(L)=L_{12}+L_{34}=\tfrac12\tr(J^{-1}L)
\]
を定める。
\begin{theorem}[一意な primitive 正規化]\label{thm:primitive}
定数項が $J$ である射影 bivector の各 perspective 軌道は、$\tau(L)=0$ の断面と
ただ一度交わる。その正規化 parameter は
\[
 a=\tfrac15\grad\tau(L)
\]
である。軌道座標 $a$ と正規化後の係数を組にした写像は、
$\AA^4$ とこの断面との直積への多項式同型である。
\end{theorem}
\begin{proof}
$(E\wedge Ja)_{12}+(E\wedge Ja)_{34}=-a^t y$ なので、
\eqref{eq:Ltransform}から
\[
 \tau(L')=\tau(L)-5a^t y
\]
を得る。これが存在と一意性を同時に示す。有限公式と逆作用 $g_{-a}$ が多項式性を与える。
\end{proof}

$[J,L]=0$ を満たす一次 bivector の空間を $Z_1$ とする。
この条件は24係数に対する独立な4線形条件なので、$\dim Z_1=20$。
独立性は $[J,-]$ を定数三重ベクトルの4係数で表す行列の階数4、または
下記の微分形式による解釈から確認できる。

\begin{corollary}\label{cor:split}
$T_a=-2(a^t y)J+E\wedge Ja$ と置くと
\[
 Z_1=\{L:[J,L]=0,\ \tau(L)=0\}\oplus\{T_a:a\in U^*\}.
\]
第一項は16次元であり、その射影は
$L\mapsto L+T_{\grad\tau(L)/5}$ である。
\end{corollary}
\begin{proof}
$T_a$ は Poisson bivector $J$ の座標変換の一次変分なので $[J,T_a]=0$。
また $\tau(T_a)=-5a^t y$ である。これにより第二項の次元は4で、
$Z_1$ 上の $\tau$ は全射となり、直和分解が従う。
\end{proof}

\subsection{2形式による記述}
symplectic lowering を行列表記で
\[
 \beta=-JLJ,\qquad L=-J\beta J
\]
と定める。$[J,L]=0$ は $\dd\beta=0$ と同値であり、
$\Lambda\beta:=\beta_{12}+\beta_{34}$ は $\tau(L)$ に等しい。
従って残る16次元空間は
\begin{equation}\label{eq:W}
 W=\{\beta\in U^*\otimes\wedge^2_0U^*: \dd\beta=0\}
\end{equation}
である。ここで $\wedge^2_0U^*$ は $\Lambda$ の核、次元5の空間である。

$AJA^t=\lambda J$ を満たす $A\in\CSp_4$ の作用を
\[
 P(y)\longmapsto\lambda^{-1}AP(A^{-1}y)A^t
\]
と正規化すると、$\tau(L)$ は $\tau(L)(A^{-1}y)$ に移る。
従って primitive gauge は $\CSp_4$ で保たれ、$W$ 上の作用は線形である。
$\Sp_4$ に限れば、2形式に対して通常の逆 pullback になる。
下記の計算証明書では pullback 作用を用いるが、逆元を取るだけなので軌道とその次元は同じである。

\begin{remark}
第2版の保存座標は、選んだ4係数を零にする別の gauge を使う。
両 gauge は上の perspective 公式で多項式的に同型となる。
しかし\eqref{eq:Qtransform}には一次項の二次式が現れるので、
保存された134式の一次変数だけを取り替える操作は正しくない。
以下では primitive gauge に移した模型を $\widetilde S_{61}$ と書く。
\end{remark}

\section{係数階数一の錐と固定 fibre}
$\beta=\sum_{i=1}^4y_i\theta_i\in W$ を線形写像
\[
 B_\beta:U\longrightarrow\wedge^2_0U^*,\qquad B_\beta(e_i)=\theta_i
\]
とみなす。係数階数 $\rank B_\beta$ は $\Sp_4$ 不変であり、0から4までを取る。
これは元の Poisson bivector の一般階数とは別の概念である。
以下の錐は被約な閉集合として扱い、determinantal ideal の根基性は主張しない。

\begin{theorem}[係数階数一の層]\label{thm:rankone}
非零 $\beta\in W$ について $\rank B_\beta=1$ であることと、
独立な一次形式 $\ell,m$ を用いて
\[
 \beta=\ell(y)\,\dd\ell\wedge\dd m,
 \qquad\langle\ell,m\rangle_J=0
\]
と書けることは同値である。この locus は5次元の単一 $\Sp_4$ 軌道である。
その射影化は、Lagrangian 二平面とその内部の直線からなる等方旗多様体
$F(1,2;4)$ に同型であり、次元は4である。
\end{theorem}
\begin{proof}
階数一なら $\beta=\ell\theta$ と書け、閉性は $\dd\ell\wedge\theta=0$ となる。
$\dd\ell$ を基底の最初に取れば、この核は $\dd\ell\wedge U^*$ なので
$\theta=\dd\ell\wedge\dd m$ と書ける。primitive 性は
$\langle\ell,m\rangle_J=0$ と同値である。逆向きは直接確認できる。

$[\beta]$ は $\langle\ell\rangle\subset\langle\ell,m\rangle$ を一意に定める。
逆に旗は $[\ell\otimes(\dd\ell\wedge\dd m)]$ を定め、これらは互いに逆の代数的写像である。
Lagrangian Grassmannian は3次元で、その上の直線の選択は $\PP^1$ だから、
旗多様体は既約4次元、対応する affine 錐は既約5次元である。
$\Sp_4$ はこの旗多様体に推移的に作用する。
代表 $\beta_0=y_1\dd y_1\wedge\dd y_3$ を保つ旗の対角群は、
その線上に指標 $(t_1,t_2)\mapsto t_1^2t_2$ で作用する。
この指標は $\CC^*$ に全射なので、非零 affine 錐自体も一軌道である。
\end{proof}

\begin{proposition}[階数一の一次項の実際の延長]\label{prop:extension}
定理\ref{thm:rankone}の各 $\beta$ に対応する $L$ について $J+L$ は Poisson であり、
その affine Pfaffian は1である。従って階数一層の全点は実際の射影 Poisson 構造へ延長する。
\end{proposition}
\begin{proof}
$X_\ell=J\grad\ell$、$X_m=J\grad m$ と置くと
$L=\ell X_\ell\wedge X_m$。
両ベクトル場は定数で可換であり、両者とも $\ell$ を消す。
従って $[L,L]=0$。閉性から $[J,L]=0$ でもある。
$L\wedge L=0$ と primitive 性から $J\wedge L=0$ なので
$(J+L)^2=J^2$、従って Pfaffian は1となる。
斉次持ち上げ $x_0^2J+x_0\ell(x)X_\ell\wedge X_m$ は射影大域切断である。
\end{proof}

具体的な代表は
\[
 L_0=y_1\partial_2\wedge\partial_4,\qquad
 \Pi_0=x_0^2(\partial_{12}+\partial_{34})+x_0x_1\partial_{24}.
\]
この延長は、同じ一次項を持つ全 Poisson 構造の分類を与えない。

\subsection{階数一層の全解を一つの fibre へ移す}
$H=\Stab_{\Sp_4}(\beta_0)$ とする。$H$ は旗の Borel 安定化群の中で、
原始指標 $t_1^2t_2$ の核である。Borel 群の unipotent radical は連結、
この指標の torus 上の核も連結なので、$H$ は連結5次元である。

\begin{corollary}[階数一層の残余問題]\label{cor:fiber}
$F$ を $\widetilde S_{61}$ において一次項が $L_0$ に等しい fibre とする。
係数階数一の層の全逆像は
\[
 \widetilde S_{61}|_{\rank B=1}\simeq\Sp_4\times^H F
\]
と表される。従ってこの層の既約成分は $F$ の既約成分と一対一に対応し、次元差は5である。
\end{corollary}
\begin{proof}
$\widetilde S_{61}\to W$ は同変写像であり、$\rank B=1$ は等質空間 $\Sp_4/H$。
同変写像の一軌道上の逆像を、その点の fibre から associated bundle として復元する。
連結な $H$ は fibre の既約成分を置換しないので、最後の主張も従う。
\end{proof}

$L_0$ は保存 gauge と primitive gauge の両方に属する。
両者を結ぶ perspective parameter は一次項だけで決まるため、この固定 fibre 全体で零である。
保存座標では $u_{34}=1$、他の15個の重み1変数を零にする。
45個の重み2変数を残して134式へ代入すると、通常次数2の134式を得る。
有理数体上の定数係数行簡約により、次を厳密に確認した\cite{fiber-certificate}。
\begin{enumerate}
\item 134式は定数係数線形独立である。
\item その線形包に非零の一次以下の多項式は存在しない。
\item 重み2変数がすべて零の既知点では Jacobian の階数は7であり、
      固定 fibre の Zariski 接空間次元は $45-7=38$ である。
\end{enumerate}
式には一次項もある。項目2は一次化の消滅を意味しない。
項目3の階数7の小行列式は $-24/25$ で、元の行列と行・列番号は証明書に保存した。
38を fibre 自身の次元や全 $S_{61}$ の接空間次元と呼ばない。
また、この層の閉包が全空間の成分を与えるかは別に判定する必要がある。

\section{Toric 成分からの一次ジェット射影の優越性}
六変数 $a,b,c,d,e,f$ に対し
\[
 K=\begin{pmatrix}
 0&a&b&c\\-a&0&d&e\\-b&-d&0&f\\-c&-e&-f&0
 \end{pmatrix},\qquad \Delta=af-be+cd
\]
と置く。開集合 $D=\{a\Delta\ne0\}\subset\AA^6$ 上で
\[
 R=\begin{pmatrix}
 1&0&0&0\\0&1/a&0&0\\d/a&-b/a&1&0\\
 e/\Delta&-c/\Delta&0&a/\Delta
 \end{pmatrix}
\]
と置けば、直接計算により $RK R^t=J$、$\det R=1/\Delta$ である。

\begin{theorem}[一般一次ジェットの toric 延長]\label{thm:toric-dominance}
既知 toric 成分 $T01$ の primitive 一次ジェット射影は $W$ 上へ優越的である。
より具体的には、明示した写像
\[
 D\times\Sp_4\longrightarrow W
\]
は優越的かつ一般に有限であり、その全点は $T01$ に属する実際の Poisson 構造から得られる。
従って $W$ のある非空 Zariski 開部分の全点は射影 Poisson 構造へ延長できる。
\end{theorem}
\begin{proof}
$u$ 座標で
\[
 P_K(u)_{ij}=K_{ij}(1+u_i)(1+u_j)
\]
を考える。これは相互可換なベクトル場 $(1+u_i)\partial_{u_i}$ の定係数外積からなる
Poisson bivector であり、斉次化の係数は $(x_0+x_i)(x_0+x_j)$ なので $\PP^4$ へ延長する。
この一般 toric 族は既知の $T01$ に含まれる。
$y=Ru$ として定数項を $J$ に移すと、$w=R^{-1}y$ を用いた一次項は
\[
 L(y)=R\bigl(\operatorname{diag}(w)K+
                 K\operatorname{diag}(w)\bigr)R^t.
\]
定理\ref{thm:primitive}を適用し、$\alpha=\frac15\grad\tau(L)$ と置いて
\[
 L_0=L-2(\alpha^t y)J+E\wedge J\alpha,\qquad\beta=-JL_0J
\]
とする。これらは $D$ 上の正則関数であり、全 parameter に対して
$\dd\beta=\Lambda\beta=0$ を満たす。有限の perspective 公式により、この一次項は実際の射影座標変換から得られる。
さらに $\Sp_4$ を作用させ、上の写像を得る。

有理点 $(a,b,c,d,e,f)=(2,3,5,7,11,13)$、$\Delta=28$ において、
六つの parameter 微分と十の $\mathfrak{sp}_4$ 作用を列に持つ24行16列行列を作る。
その16次小行列式は厳密に
\[
 \frac{258241}{10504375}\ne0
\]
である。元の小行列、行・列番号、全 parameter での写像と制約の確認を
\cite{toric-certificate}に保存した。
像は16次元の $W$ に含まれ、その点で微分が全射なので像の閉包は $W$ である。
定義域と値域はともに既約16次元だから写像は一般に有限。
構成可能像の定理により、像は非空 Zariski 開部分を含む。
\end{proof}

\begin{remark}
既存の成分対応と $\dim T01=25$ を併用すると、対応する16次元の $S_{61}$ 成分自身も
$W$ に優越的かつ一般に有限に写る。
一般次数は後の定理\ref{thm:toric-degree}で3と決定する。
分岐と、像に含まれる開部分の具体的な方程式は未計算である。
任意の一次ジェットの延長、延長の一意性、他成分由来の延長の不存在は主張しない。
特に「一般の一次項は toric 成分で実現する」ことから、
「全 Poisson 構造が toric 成分に属する」とは従わない。
\end{remark}

\section{Toric 開部分の16変数 fibre}\label{sec:coframe}
一次部分 $\beta\in W$ を固定する。この節で扱うのは、実際の退化因子が
射影標構をなす五つの超平面の和であり、基点0を避ける構造である。
$T01$ のこの稠密開部分に対象を限り、全 $S_{61}$ の45変数 fibre を同じ問題へ
帰着したとは主張しない。

\subsection{中心化した五超平面と全双ベクトルの復元}
各超平面の方程式を一意に $h_i=1+\ell_i(y)$、$0\le i\le4$ と規格化する。
$P(0)=J$ の場合、その affine Pfaffian は $\prod_i h_i$ である。
一次項を比較すると primitive 条件は $\sum_i\ell_i=0$ と同値となる。
$A$ の第$i$行を $a_i=\ell_i-\ell_0$、$1\le i\le4$ とする。
射影標構の条件は $\det A\ne0$ であり、逆に $A\in\GL_4$ から
\begin{equation}\label{eq:coframe-h}
 z=Ay,\qquad h_0=1-\frac15\sum_{i=1}^4z_i,\qquad h_i=h_0+z_i
\end{equation}
によって中心化した五超平面を復元できる。

$\mathbf1^t=(1,1,1,1)$ と置く。
\begin{proposition}[toric 開部分の復元]\label{prop:coframe}
各 $A\in\GL_4$ に対して
\begin{align}
 F_A(y)&=A^{-1}\left(I-\frac15z\mathbf1^t\right)
          \operatorname{diag}(h_1,\ldots,h_4)A,\nonumber\\
 P_A(y)&=F_A(y)J F_A(y)^t\label{eq:coframe-P}
\end{align}
は射影 Poisson 構造を与え、定数項は $J$、一次項は primitive、
実際の退化因子は $\prod_i h_i=0$ である。
逆にこの五超平面配置と定数項 $J$ を持つ Poisson 構造は一意であり、
\eqref{eq:coframe-P}に等しい。
\end{proposition}
\begin{proof}
標準 toric 座標を $t_i=h_i/h_0$ とする。
$\Lambda_A=AJA^t$ として $\{t_i,t_j\}=(\Lambda_A)_{ij}t_it_j$ と置くと、
相互可換な $t_i\partial_{t_i}$ により Jacobi 恒等式を満たす。
$u_i=t_i-1=z_i/h_0$ の微分を使って $y$ 座標へ戻すと\eqref{eq:coframe-P}を得る。
これは射影座標変換であるから大域に延長する。
$F_A$ の二次部分は階数一の外積であり、$P_A$ の四次部分は零、従って次数は高々3である。
行列式補題により
\[
 \det F_A=h_0\prod_{i=1}^4h_i,\qquad
 \operatorname{Pf}(P_A)=\det(F_A)\operatorname{Pf}(J)=\prod_{i=0}^4h_i.
\]
各因子は相異なり、一次項の和は零である。

逆に Hamiltonian 場は $P$ と $P\wedge P$ を保存するので、
その退化因子の各既約成分を保存する。
従って被約な五超平面配置に対して $P$ は対数的な双ベクトルである。
標準 toric 境界の対数接束は $t_i\partial_{t_i}$ によって自明化されるので、
大域対数双ベクトルはこれらの定係数外積の和である。
基点での値 $J$ が六つの係数を $\Lambda_A=AJA^t$ と一意に決める。
\end{proof}

\subsection{固定一次項の方程式とラベルの除去}
\[
 H(z)=\operatorname{diag}(z)-\frac15(\mathbf1^tz)I
                         -\frac15z\mathbf1^t,\qquad C_A(y)=A^{-1}H(Ay)A
\]
と置くと、\eqref{eq:coframe-P}の一次項は $C_AJ+JC_A^t$、対応する2形式は
\begin{equation}\label{eq:coframe-beta}
 \beta_A=JC_A+C_A^tJ
\end{equation}
である。$\operatorname{tr}H=0$ と Jacobi 恒等式から $\beta_A\in W$ となる。
従って固定した $\beta$ の coframe fibre は
\begin{equation}\label{eq:coframe-fiber}
 J\operatorname{adj}(A)H(Ay)A+
   \bigl(\operatorname{adj}(A)H(Ay)A\bigr)^tJ
       -\det(A)\beta(y)=0,\qquad \det A\ne0
\end{equation}
で与えられる。六つの行列成分の四つの一次係数を比較した24式には、
閉性と primitive 性の8本の固定線形関係がある。
$W$ の線形座標を選べば16式となり、未知数は $A$ の16成分、通常次数は5次と4次である。
$\det A$ による局所化を保持し、特異行列の解を延長と数えない。

\begin{theorem}[toric fibre と有限対称群]\label{thm:coframe-fiber}
\eqref{eq:coframe-fiber}の解集合の自由な $S_5$ 作用による軌道は、
対象とした toric 開部分の、一次項 $\beta$ を持つ延長と一対一に対応する。
$\GL_4\to W$、$A\mapsto\beta_A$ は優越的かつ一般に有限である。
その一般次数を $d$ とすると、$T01$ の一次部分射影の一般次数は $d/120$ である。
\end{theorem}
\begin{proof}
$a_0=0$ とすると、五超平面の再ラベル $\sigma$ は
$a'_i=a_{\sigma(i)}-a_{\sigma(0)}$、すなわち $A'=M_\sigma A$ と表される。
$M_\sigma$ は標準4次元表現の整数行列で、この規約では右作用となる。
$A$ が可逆なので非自明な置換は固定点を持たない。
同じ五超平面と定数項 $J$ からの復元は一意であるから、この作用は全 $P_A$ を保つ。
逆に同じ $P_A$ を与える二つの表示は、実際の退化因子の相異なる既約因子から
五超平面を復元すると、再ラベルだけ異なる。これで全単射を得る。

定理\ref{thm:toric-dominance}の有理点における Darboux frame の逆行列を $A$ とすると、
\eqref{eq:coframe-beta}を16個の行列成分で微分した行列の16次小行列式は
\[
 \frac{258241}{294122500}\ne0
\]
となる\cite{coframe-certificate}。従って写像は優越的かつ一般に有限である。
標数零で一般 fibre は被約であり、自由な再ラベルの軌道は各120点を持つ。
$T01$ の toric 開部分の補集合は次元15以下で、一般一次項上から除けるので、
次数の関係も従う。
\end{proof}

一般の $\Lambda_A$ に対する Jacobi と Pfaffian の恒等式、上の小行列式、
隣接交換による全 $P_A$ の不変性の有理点での別検算を\cite{coframe-certificate}に保存した。
一般次数は定理\ref{thm:toric-degree}により $d=360$ となる。
120個のラベル違いを120個の異なる延長と数えない。

\section{ラベルを持たない20変数の結合代数模型}\label{sec:algebra}
前節のcoframeを、五超平面を順序付けない可換代数の積へ置き換える。
この変更により、固定一次部分のtoric延長は20変数・96二次式の開部分で表される。
以下は\cite{algebra-certificate,algebra-review}の構成と独立点検を統合したものである。

\subsection{一次写像の右逆と対称三階テンソル}
$V=\CC^4$ とし、可換双線形写像 $m:\operatorname{Sym}^2V\to V$ に対して
$m_u(v)=m(u,v)$、$m_i=m_{e_i}$ と書く。
$\tr m_u=0$ の四条件は独立であり、この空間の次元は36である。
実際、一次形式 $a$ に対して $m(u,v)=(a(u)v+a(v)u)/5$ とすれば $\tr m_u=a(u)$ となる。

\begin{lemma}\label{lem:algebra-lift}
写像 $m\mapsto\beta_m$ を
\[
 \beta_m(y)=Jm_y+m_y^tJ
\]
で定めると、trace-freeな可換積の空間から $W$ への全射となり、
核は $\operatorname{Sym}^3V^*$、次元20である。
$\beta_{jk,i}$ を $\beta_{jk}(y)$ の $y_i$ 係数とすれば、全持ち上げは
\begin{align}
 T_{ijk}&=\frac{\beta_{jk,i}+\beta_{ik,j}}3+c_{ijk},\qquad
 c_{ijk}=c_{\sigma(i)\sigma(j)\sigma(k)},\label{eq:algebra-lift}\\
 m(e_i,e_j)&=J(T_{ij1},\ldots,T_{ij4})^t\nonumber
\end{align}
で与えられる。
\end{lemma}
\begin{proof}
$T_{ijk}=m(e_i,e_j)^tJe_k$ と置くと、可換性は $T_{ijk}=T_{jik}$、
$\beta_{jk,i}=T_{ijk}-T_{ikj}$ である。
この式の巡回和は零なので $\dd\beta=0$ であり、成分を比較すると
$\beta_{12,i}+\beta_{34,i}=\tr m_i$ を得る。
閉性を使えば\eqref{eq:algebra-lift}の $c=0$ 部分が右逆となる。
核では $T$ が第一・第二添字と第二・第三添字の交換に不変だから、完全対称である。
独立係数の個数は $\binom63=20$。primitive条件からtrace-free性も従う。
\end{proof}

\subsection{結合性を表す96二次式}
固定した $\beta$ について、\eqref{eq:algebra-lift}の20係数 $c_{ijk}$ を変数とする。
\begin{equation}\label{eq:algebra-quadrics}
 q_{ij}=\tfrac13\tr(m_i m_j),\qquad
 C_{ij}=[m_i,m_j]+e_iq_j-e_jq_i=0\quad(i<j)
\end{equation}
と置く。$q_j$ は対称行列 $q$ の第$j$行である。
六つの4次行列の全成分を取るため96本の二次式であり、96本の独立性は主張しない。

\begin{proposition}\label{prop:associative}
\eqref{eq:algebra-quadrics}は、$B=\CC1\oplus V$ 上の可換単位付き積
\begin{equation}\label{eq:algebra-product}
 (a,u)(b,v)=(ab+q(u,v),\ av+bu+m(u,v))
\end{equation}
の結合性と同値である。省略した結合性条件も、96式の生成イデアルに属する。
\end{proposition}
\begin{proof}
$e_i$ による左乗法の行列を
\[
 L_i=\begin{pmatrix}0&q_i\\ e_i&m_i\end{pmatrix}
\]
とする。$[L_i,L_j]$ の左上と左下はqの対称性、mの可換性により零で、
右下は $C_{ij}$、右上は $D_{ij}=q_i m_j-q_jm_i$ である。
一般の可換mと対称qに対して、行列Jacobi恒等式と $m_i e_j=m_j e_i$ から
\[
 e_iD_{jk}+e_jD_{ki}+e_kD_{ij}
 =-\bigl([m_i,C_{jk}]+[m_j,C_{ki}]+[m_k,C_{ij}]\bigr)
\]
を得る。相異なる $i,j,k$ の第$i$行を取れば、$D_{jk}$ の各成分は $C$ の成分の
生成イデアルに属する。従って全 $L_i$ は可換となる。
可換積では左乗法の可換性から $x(yz)=y(xz)=(xy)z$ が従い、結合的である。

逆に結合的なら $C_{ij}=0$。
trace-free性から正則表現のトレースは $\operatorname{Tr}L_{(a,u)}=5a$ であり、
結合性を用いると $\operatorname{Tr}(L_iL_j)=5q_{ij}$。
ブロックを掛けて計算した値は $2q_{ij}+\tr(m_i m_j)$ だから、qの式も必然的に成立する。
\end{proof}

\begin{theorem}[toric延長の20変数模型]\label{thm:algebra-model}
\eqref{eq:algebra-quadrics}の $\det q\ne0$ の開部分は、固定一次部分 $\beta$ を持つ
前節のtoric延長のスキームを与える。120通りの再ラベルは既に除かれている。
全双ベクトルは文字の計算を使わず
\begin{equation}\label{eq:algebra-P}
 F_m(y)=I+m_y-yq_y,\qquad P_m(y)=F_m(y)JF_m(y)^t
\end{equation}
で復元でき、そのPfaffianは $\operatorname{Norm}_B(1+y)$ である。
\end{theorem}
\begin{proof}
正則表現のトレースペアリングは $(1,V)$ の基底で $\operatorname{diag}(5,5q)$ となる。
$\det q\ne0$ なら非退化。冪零元 $n$ に対し、任意の $b$ の $nb$ も冪零で
$\operatorname{Tr}L_{nb}=0$ だから、非退化性により非零冪零元は存在しない。
従って $B\simeq\CC^5$。
五つの文字を $\chi_i$、そのVへの制限を $\ell_i$ とすれば
$\sum_i\ell_i=0$ で、Vは $\CC^5$ のトレース零部分へ同型に写る。
$h_i=1+\ell_i$ として五超平面を回復し、行 $\ell_i-\ell_0$ によるcoframe Aは可逆となる。
逆にcoframeから $\CC^5$ の成分ごとの積を移せば、同じmを回復する。
文字の順序を変えてもmは変わらない。

判別式条件は
\[
 q=A^t\frac{5I-\mathbf1\mathbf1^t}{25}A,\qquad
 \det q=\frac{(\det A)^2}{5^5}
\]
となり、前節の開条件に一致する。
$L_i$ が $\PP(B)$ のchart $[1:y]$ に誘導する射影ベクトル場は
$e_i+m_i y-yq_i(y)$。これらは可換で、その列行列が $F_m$ なので
\eqref{eq:algebra-P}は射影Poisson構造である。
Schur補元により
\[
 \det F_m=\det\begin{pmatrix}1&q_y\\y&I+m_y\end{pmatrix}
           =\operatorname{Norm}_B(1+y)=\prod_i(1+\ell_i(y)).
\]
この実際の退化因子から文字の無順序集合を復元でき、mも一意に決まる。

点の全単射だけではスキームとしての主張は足りないので、微小変形も確認する。
剰余体 $\CC$ の局所Artin代数R上で $\det q$ が単元である解を考える。
特殊ファイバーの五つの原始冪等元は、$e^2-e$ の導関数 $2e-1$ が単元なので
一意に持ち上がり、直交性と和1を保つ。各直和因子はR上自由階数1であり、
持ち上げた冪等元を基底に取ると代数全体は $R^5$ となる。
正則トレースの核は $V\otimes R$ であり、文字から可逆coframeを回復できる。
逆の構成も自然で、順序だけが自由な $S_5$ 作用に対応する。
従って半単純開部分の変形関手と完成局所環はcoframe商に一致し、
一次部分写像も一致する。特に、βを動かしたこの開模型は滑らかな16次元空間である。
\end{proof}

\begin{remark}
\eqref{eq:algebra-P}は $\det q=0$ の結合代数からも射影Poisson構造を作るが、
その部分の全成分、toric境界への所属、全S61を尽くすかは別問題である。
非零判別式上のスキーム同定を、判別式零の全模型へ無条件に延長しない。
\end{remark}

\section{Toric成分の一般次数は三}\label{sec:toric-degree}
20変数模型の次数を、特殊ファイバーの全イデアルと、その近傍での有限性から決定する。
一つの特殊ファイバーの点数だけでは、無限遠へ逸走する解を除外できない。
以下ではその点を含む開集合上で\emph{全解スキーム}が有限となることを先に示す。

\begin{theorem}[toric一次射影の一般次数]\label{thm:toric-degree}
$T01$ に対応する成分から $W$ への一次部分射影の一般次数は3である。
従って十分一般の一次部分は、ちょうど三つの相異なるtoric延長を持つ。
順序付きcoframe写像 $\GL_4\to W$ の一般次数は360である。
\end{theorem}
\begin{proof}
\textbf{普遍的な一次消去。}
\eqref{eq:algebra-lift}の20係数を辞書順に $c_0,\ldots,c_{19}$ と並べる。
各乗法行列は $m_i=m_i^0(\beta)+m_i^1(c)$ と線形に分かれるので、96式のcに関する
二次斉次部分はβに依存しない。
旧toric係数 $(2,3,5,7,11,13)$ に対応する一次部分を $\beta_*$ とする。
この点で96式の定数係数線形包の階数は63、二次斉次部分の階数は45であり、
18本の独立な一次帰結を得る。各帰結を96式の有理定数係数線形結合として表した
18組の係数を\cite{degree-certificate}に保存した。

二次部分のβ非依存性により、\emph{同じ}18組の線形結合は全βで二次項を消す。
従って普遍的一次帰結 $M(\beta)c+b(\beta)=0$ を与える。
$M$ の成分はβの一次式、bの成分は二次式である。
$c_0,\ldots,c_{17}$ の18次小行列式を $\delta(\beta)$ とすると、保存した行順で
$\delta(\beta_*)=-1$ である。$\delta\ne0$ に局所化し、これら18変数を
$x=c_{18},y=c_{19}$ のアフィン一次式として消去する。
この消去は元のイデアルの帰結を用い、残る全96式も保持するので、解を落とさない。

\textbf{全解スキームの有限性。}
消去後の96式はx,yの高々二次式で、係数はβの正則有理関数である。
$\beta_*$ において、その最初の三式の $(x^2,xy,y^2)$ の係数行列は
\[
 N_*=
 \begin{pmatrix}
 73/8232&43/43218&5717/2420208\\
 -131/4116&164/21609&-1763/1210104\\
 109/8232&-1/1764&47/2420208
 \end{pmatrix},\qquad
 \det N_*=-\frac{41639}{189823367328}\ne0.
\]
同じ三式の係数行列 $N(\beta)$ の行列式 $\epsilon$ も非零となる開集合へ縮める。
すると $x^2,xy,y^2$ は $1,x,y$ の線形結合として表される。
次数に関する帰納法から、全96式による商代数は基礎の局所化座標環上で
$1,x,y$ により生成される有限加群となる。
よって $\beta_*$ を含む開集合上の二次模型の全解スキームは有限である。
未知の追加の一次帰結が存在しても、この有限性を妨げない。

\textbf{特殊ファイバーの全解。}
$\beta_*$ で上の三式を正規化すると
\begin{align*}
 x^2&=\tfrac{119}{39}x-\tfrac3{13}y+\tfrac{4214}{117},\\
 xy&=-\tfrac{490}{13}x-\tfrac{119}{39}y-\tfrac{6860}{39},\\
 y^2&=\tfrac{3430}{13}x+\tfrac{490}{13}y+\tfrac{48020}{39}
\end{align*}
となる。全96残余式がこの三式の定数係数線形結合であることを検算し、
得られたx,yの乗法行列の可換性も確認した。双方のイデアル包含により、全残余イデアルは
\begin{equation}\label{eq:toric-sample-ideal}
 \left(x-\frac{13}{3430}y^2+\frac17y+\frac{14}3,
       \quad y^3-\frac{24010}{13}y\right)
\end{equation}
に等しい。商代数の基底は $1,x,y$ で、yは巡回元である。
三次式 $g(y)=y^3-24010y/13$ の判別式は
\[
 \operatorname{disc}(g)=\frac{55365148804000}{2197}\ne0.
\]
従って全特殊ファイバーは三つの相異なる被約複素点からなる。
既知点は $(x,y)=(-14/3,0)$、他の二点は $y=\pm\sqrt{24010/13}$ と
\eqref{eq:toric-sample-ideal}で与えられる。

さらに $\det q$ のgによる剰余は
\[
 r(y)=\frac{11713280438427y^2-513182409985385y+21694099123200}
                {86472015000000}.
\]
有理係数B\'ezout恒等式によって $\gcd(r,g)=1$ を検算した。
従って三点の全てが $\det q\ne0$ の半単純部分に属し、特殊ファイバーに捨てる点はない。

\textbf{有限\'etaleな近傍への移行。}
上の有限射を $X\to B$ とし、Bは $W$ の $\beta_*$ を含む非空開集合とする。
$V(\det q)\subset X$ の像は有限射のため閉であり、$\beta_*$ を含まないのでBから除ける。
残る始域は、定理\ref{thm:algebra-model}によりcoframe商と同じ局所スキーム構造を持つ
滑らかな16次元スキームである。
$\beta_*$ 上の三点では、被約な零次元ファイバーの接空間は零である。
始域と終域がともに滑らかな16次元なので、射の微分は同型で、各点で\'etaleとなる。
非\'etale部分も閉であり、その有限像は $\beta_*$ を避けるので除ける。
こうして $\beta_*$ を含む開集合上で射は有限\'etaleとなる。

非空のW開集合は既約・連結であるから、この有限\'etale射の次数は一定であり、
$\beta_*$ 上の三点からその値は3となる。
toric開部分はT01で稠密で、その補集合は一般一次部分上から除けるので、
これはT01成分自身の一次射影の一般次数である。
自由な $S_5$ 被覆を戻すと、順序付きcoframeの次数は $120\cdot3=360$ となる。
\end{proof}

定数行操作を追跡した証明書\cite{degree-certificate}に加えて、
別の実装\cite{degree-review}では二次斉次係数行列の左核から一次帰結を構成し、
同じ18変数消去、全残余イデアル、三次式の分離性、$\det q$ との互いに素性を確認した。
この結果はtoricの一般次数を決めるが、分岐因子、例外ファイバー、
二つの楕円型成分の次数、全S61ファイバーの尽くしは決めない。

\section{二つの楕円型成分の優越性}\label{sec:elliptic}
$T39$ と $T40$ は内部カタログの名前であり、それぞれ $q_{5,1}$、$q_{5,2}$ の成分を表す。
使用する斉次係数族は Ortenzi--Rubtsov--Tagne Pelap\cite{ortp}の(16)、(18)、(19)による。
添字を $\mathbb Z/5\mathbb Z$ とし、反対称行列 $B$ を
\begin{align}
 B_{i,i+1}&=a_1x_ix_{i+1}+a_2x_{i+2}x_{i+4}+a_3x_{i+3}^2,\nonumber\\
 B_{i,i+2}&=b_1x_ix_{i+2}+b_2x_{i+3}x_{i+4}+b_3x_{i+1}^2\label{eq:elliptic-family}
\end{align}
と定める。$\lambda\ne0$ に対し係数は次表で与えられる。
\begin{center}
\small
\begin{tabular}{@{}c*{6}{c}@{}}
\toprule
族 & $a_1$ & $a_2$ & $a_3$ & $b_1$ & $b_2$ & $b_3$\\\midrule
$q_{5,1}$ & $-\frac35\lambda^2+\frac1{5\lambda^3}$ & $-2/\lambda$ & $1/\lambda^2$
& $-\frac15\lambda^2-\frac3{5\lambda^3}$ & $2$ & $\lambda$\\[2mm]
$q_{5,2}$ & $\frac25\lambda^2+\frac1{5\lambda^3}$ & $\lambda$ & $-1/\lambda$
& $-\frac15\lambda^2+\frac2{5\lambda^3}$ & $-1/\lambda^2$ & $1$\\
\bottomrule
\end{tabular}
\end{center}

\begin{remark}[一次資料の符号]
\cite{ortp}の(18)では $q_{5,1}$ の $b_3$ は $+\lambda$ であるが、
後の展開式(21)では $-\lambda$ と印刷されている。本稿は(18)を採用し、
10個の斉次 Jacobi 式を $\QQ(\lambda)[x_0,\ldots,x_4]$ 上で恒等的に検証した。
$-\lambda$ へ替えると $(0,1,2)$ 成分は
\[
 -4(\lambda x_0^2x_3-\lambda x_2^2x_4+x_0x_4^2-x_2x_3^2)
\]
となり非零である。この検算と符号の比較も\cite{elliptic-certificate}に含む。
\end{remark}

$x_0=1$ の座標を $z_i=x_i/x_0$ とすると、降下した双ベクトルは
\[
 P_{ij}(z)=B_{ij}(1,z)-z_iB_{0j}(1,z)+z_jB_{0i}(1,z)
 \quad(1\le i,j\le4)
\]
であり、四つの affine Jacobi 式も一般 $\lambda$ で検証できる。
基点 $p\in\AA^4$、非零スカラー $c$ を動かし、$K=cP(p)$ が非退化な開部分で
$FKF^t=J$ を満たす有理 frame を選ぶ。
具体的には座標順を $(1,4,2,3)$ に替える置換行列 $R_0$ を用い、
$R_0KR_0^t$ の上三角を $(a,b,c',d,e,f)$、$\Delta=af-be+c'd$ と書けば
\[
 F=\begin{pmatrix}
 1&0&0&0\\0&1/a&0&0\\d/a&-b/a&1&0\\
 e/\Delta&-c'/\Delta&0&a/\Delta
 \end{pmatrix}R_0
\]
が $a\Delta\ne0$ 上で使える。$y=F(z-p)$ における一次項の係数行列は
\[
 L_j=F\left(\sum_h(F^{-1})_{hj}\,c\,\partial_hP(p)\right)F^t.
\]
これを定理\ref{thm:primitive}で正規化し、$\beta=-JL_{\rm prim}J$ とする。
残存 $\Sp_4$ 作用を加えると、非空開集合
\[
 \mathcal V\subset\GG_{m,\lambda}\times\AA^4_p\times\GG_{m,c}\times\Sp_4
 \longrightarrow W
\]
への正則写像を得る。$c$ は新しい楕円曲線のモジュライではなく、
正規化後の座標拡大を動かすパラメータである。
全操作が実際の Poisson 構造の射影座標変換と正規化なので、
その像は対応する既知成分の一次像に含まれる。

\begin{theorem}[楕円型成分の一般有限な一次射影]\label{thm:elliptic-dominance}
$T39,T40$ に対応する16次元の $\widetilde S_{61}$ 成分は、ともに $W$ に優越的かつ
一般に有限に写る。
\end{theorem}
\begin{proof}
両族とも $\lambda=2$、$p=0$、$c=1$ を用いる。
この点では $K_{14}=-b_3$、$K_{23}=a_3$、他の上三角成分は零であり、
$K$ は非退化で上の frame の開条件を満たす。
六つのパラメータ $(\lambda,p_1,p_2,p_3,p_4,c)$ の微分と
十個の $\mathfrak{sp}_4$ 作用を列に持つ24行16列行列を作る。
二形式の対を辞書順、その中で一次変数を $y_1,\ldots,y_4$ 順に並べ、
零始まりの行番号
\[
 (0,1,2,3,4,5,6,7,8,9,11,13,14,15,17,19)
\]
を選ぶと、小行列式はそれぞれ
\[
 T39:\quad\frac{142978515625}{512},\qquad
 T40:\quad-\frac{142978515625}{33554432}
\]
である\cite{elliptic-certificate}。どちらも非零なので微分は $W$ へ全射であり、
像は稠密である。既知の射影成分次元25と次元差9により、各 slice 成分の次元は16。
従って優越写像は一般に有限である。
\end{proof}

frame 微分も証明書の一部である。$H=(\delta F)F^{-1}$ に対し
$H_0=\frac12F(\delta K)F^tJ$ を使うと
\[
 H_0J+JH_0^t+F(\delta K)F^t=0,\qquad H-H_0\in\mathfrak{sp}_4
\]
となる。$H$ と $H_0$ による一次部分微分の差が、既に含めた $\Sp_4$ 作用列の
線形結合になることを直接検算し、両方法の16次小行列式の一致も確認した。

\section{全 $S_{61}$ スキーム内の滑らかな非分岐点}\label{sec:etale}
ここでは族のパラメータ微分に加え、保存された全134式を61変数で微分する。
$T01$ の係数 $(2,3,5,7,11,13)$ と、前節の $T39,T40$ の $\lambda=2$、$p=0$ の点を、
それぞれ定数項 $J$ に正規化して旧 S61 gauge へ移す。

\begin{theorem}[全スキームでの正則性と一次射影]\label{thm:etale}
上記の三点は全 $S_{61}$ の滑らかな点で、一次部分への射影は各点で \'etale である。
従ってその fibre における各点は孤立した被約点である。
この性質は各既知成分の非空開部分で成り立つ。
\end{theorem}
\begin{proof}
全134式の61変数 Jacobian は各点で階数45である。
重み2の45列だけに制限しても階数45となる。
この後者の45次小行列式は、保存した行順序に対して次の値である。
\begingroup\small
\begin{align*}
 T01:&\quad-\frac{4696557768725188309033758541443977502605228011369260075725936199}
                       {7444176447590624000000},\\
 T39:&\quad-\frac{14489624420239116609310651127244682356831617653369903564453125}
                       {75557863725914323419136},\\
 T40:&\quad-\frac{142207129520970493638041972637771701965903048403561115264892578125}
                       {1329227995784915872903807060280344576}.
\end{align*}
\endgroup
有理数体上の全行列、行・列番号、各点の61座標は\cite{known-fiber-certificate}にある。
同じ証明書は134式の消滅、55復元式による150係数の一致、
独立な projective Jacobi の消滅も確認している。

従って全スキームの接空間次元は $61-45=16$ である。
他方、各点は既知の16次元 slice 成分に属するので、局所次元は16以上。
局所次元と埋込み次元がともに16となり、その局所環は正則である。
$\CC$ 上の有限型スキームなので、各点は滑らかであり、
別の成分や非被約構造はその点を通らない。

旧 gauge の16個の重み1係数と primitive 一次部分は線形同型で結ばれる。
重み2の45列の階数45は、一次部分を固定した fibre の接空間が零であることを表す。
従って16次元の滑らかな始域と $W$ の間の微分は同型であり、射影は \'etale となる。
滑らかさと \'etale 性は開条件であるから、最後の主張も従う。
\end{proof}

この証明は局所次元の下界として既知成分への所属を用いる。
接空間計算だけから成分性を導いたものではない。
また、一点の fibre 接空間が零でも、同じ fibre の他の点や未知成分上の正次元部分は除外されない。

\section{一般一次ジェットの残存群軌道}
\subsection{一般一次ジェットでは有限個の軌道に還元できない}
$W$ の係数を、添字対 $(12),(13),(14),(23),(24),(34)$ の順に、
各対の中では $y_1,y_2,y_3,y_4$ の順に並べる。
次の一点 $\beta_*$ は $W$ に属する。
\[
\begin{split}
 (-13,-14,-15,-16;&\ 1,-9,2,-9;\ 3,-7,4,5;\\
 &6,7,8,-3;\ 9,10,11,12;\ 13,14,15,16).
\end{split}
\]
\begin{proposition}\label{prop:generic-orbit}
$W$ の稠密開集合上で $\Sp_4$ 軌道の次元は10である。
さらにスカラー作用を加えると一般軌道の次元は11である。
従って一般軌道の余次元はそれぞれ6、5である。
\end{proposition}
\begin{proof}
$\mathfrak{sp}_4$ の10個の基底として、対称行列の標準基底 $S$ に対する $JS$ を用いる。
一次2形式 $\beta=\sum y_i\beta_i$ の pullback の無限小作用は
\[
 (A\cdot\beta)_k=A^t\beta_k+\beta_k A+\sum_i A_{ik}\beta_i.
\]
$\beta_*$ における24行10列の作用行列の10次小行列式として
$-1500934987656\ne0$ を得る。
スカラー方向 $\beta_*$ を追加した行列の11次小行列式は
$-33302699006080\ne0$。
群の次元がそれぞれ10、11なので、この階数は上限に達している。
行列式の非消滅は稠密開条件であり、主張が従う。
元の行列・添字・有理数体上の計算は\cite{weight1-certificate}に記録した。
\end{proof}

\begin{remark}[延長可能性の限定]
$\beta_*$ の係数写像の階数は4で、その4次小行列式は $-1260$ である。
この特定の $\beta_*$ が実際の Poisson 構造へ延長するかは、上の計算だけでは判定していない。
一方、定理\ref{thm:toric-dominance}により実際の一次ジェット像は非空開部分を含むので、
命題\ref{prop:generic-orbit}の一般軌道次元は延長可能な一次ジェットの稠密部分でも実現する。
実際に同定理の toric 有理点でも、係数階数4、$\Sp_4$ 作用階数10、スカラー込みの階数11を
確認した。これらを全 Poisson 構造自身の軌道次元やモジュライ次元として引用してはいけない。
ここでの余次元は $W$ と軌道の次元差であり、具体的な商や不変式環の構成ではない。
\end{remark}

\section{二つの既知28次元成分の共通一次像}\label{sec:ambient}
この節では、第6版\cite{ambient-source}における既知の全成分
$\mathcal C_A,\mathcal C_B$ を扱う。成分性、相異性、射影次元28は既存の証明を入力とする。
新たに決定するのは、対応する $S_{61}$ 成分の primitive 一次ジェット像である。

\subsection{原族と共通の六次元線形部分空間}
斉次座標を $(x_0,x,y,z,w)$ とする。
$f=f(x_0,x)$ は二次、$h=h(x_0,x,z,w)$ は三次斉次式とし、
$X_h=-h_w\partial_z+h_z\partial_w$、$E_{zw}=z\partial_z+w\partial_w$ と置く。
原族は
\begin{align}
 P_A&=f\partial_x\wedge\partial_y+x_0^2\partial_z\wedge\partial_w
                  +\partial_y\wedge X_h,\label{eq:CA}\\
 P_B&=f(\partial_x\wedge\partial_y+\partial_z\wedge\partial_w)
           +\partial_y\wedge\left(-\tfrac12f_xE_{zw}+X_h\right).\label{eq:CB}
\end{align}
両者の Jacobi 恒等式は、$\operatorname{div}_{zw}X_h=0$ と
$\operatorname{div}_{zw}(-f_xE_{zw}/2)=-f_x$ から直接従う。
その実際の退化五次式はそれぞれ $x_0^3f$、$x_0f^2$ である。

$x_0=1$、$f=1+px+O(x^2)$ とし、$h$ の一次微分が原点で零となる場合をまず考える。
一次ジェットに寄与する二次部分を
\[
 h_2=A xz+B xw+Cz^2+Dzw+Fw^2
\]
とする。$x^2$ 項は $X_h$ に寄与しない。
primitive 正規化後の一次項は両族とも
\begin{equation}\label{eq:common-plane}
 L_0=tx(\partial_{xy}-\partial_{zw})
       +\partial_y\wedge\left(X_{h_2}+\tfrac t2E_{zw}\right)
\end{equation}
の形になる。$A$ 族では $p=5t/2$、$B$ 族では $p=-5t$ である。
これは第3節の有限公式を代入した全 parameter の恒等式である。

$\beta=-JL_0J$ の六次元線形空間を $S\subset W$ とする。
その成分は
\begin{gather}
 \beta_{12}=tx,\quad \beta_{34}=-tx,\quad
 \beta_{23}=\beta_{24}=0,\label{eq:S-form}\\
 \beta_{13}=(h_2)_z+\tfrac t2w,\qquad
 \beta_{14}=(h_2)_w-\tfrac t2z.\nonumber
\end{gather}
各 $t,h_2$ は、原族で $f=x_0^2+px_0x$、$h=x_0h_2$ とすれば実現する。
原点での値は常に $J$ であり、rank4 の開条件を失わない。

\subsection{任意の rank4 点における正規化}
原族の rank4 点では必ず $x_0\ne0$ である。
affine translation でその点を原点へ移しても、$A$ 族の形は保たれる。
$B$ 族では $E_{zw}$ の平行移動により生じる定数ベクトル項を $X_h$ に吸収できる。
具体的には $z,w$ の定数移動による項は $f_x(x)$ 倍の定数ベクトルであり、
$x$ に一次以下、$z,w$ に一次の Hamiltonian を加えればよい。
従って $h$ の三次以下という許容範囲を保つ。

原点での値を
\[
 K=\begin{pmatrix}
 0&a&0&0\\-a&0&d&e\\0&-d&0&g\\0&-e&-g&0
 \end{pmatrix},\qquad ag\ne0
\]
とする。$A$ 族では $g=1$、$B$ 族では $g=a$。
定数交差項 $d,e$ を零と仮定せず、
\[
 R=\begin{pmatrix}
 1&0&0&0\\0&1/a&0&0\\d/a&0&1&0\\e/(ag)&0&0&1/g
 \end{pmatrix}
\]
で正規化すると $RKR^t=J$ を得る。
この変換は $x$ と $y$ への非依存性を保つので、その後の一次項は
\[
 L_{12}=sx,\quad L_{34}=rx,\quad L_{13}=L_{14}=0,
 \quad(L_{23},L_{24})=Y(x,z,w)
\]
と書ける。Jacobi の定数項は $\operatorname{div}_{zw}Y=-r$ を与える。
$\tau=(s+r)x$ による primitive 正規化を行うと
\[
 t=\frac{2s-3r}{5},\qquad
 (L_0)_{12}=tx,\quad(L_0)_{34}=-tx,\qquad
 \operatorname{div}_{zw}Y_0=t.
\]
従って $Y_0-tE_{zw}/2$ は divergence-free な線形場であり、
二次式 $h_2$ の $X_{h_2}$ として表される。
これで任意の rank4 点からも\eqref{eq:common-plane}に到達する。
一般の $a,d,e,g$ に対するこの計算は\cite{ambient-certificate}で記号的に確認した。

同じ点を $p$ に移して定数項を $J$ にする二つの射影座標系の差は、
$p$ の安定化群に属する。perspective 部分は primitive gauge で一意に除かれ、
残る作用は $\CSp_4$ である。
$S$ は線形部分空間でスカラー倍を許すため、その $\Sp_4$ 飽和と $\CSp_4$ 飽和は一致する。

\subsection{閉 incidence と正確な像}
定数2形式を $\omega=\dd x\wedge\dd y+\dd z\wedge\dd w$ とする。
$v\in U\setminus\{0\}$ に対し $\ell_v=\iota_v\omega$ と置き、
\[
 S_{[v]}=\{\beta\in W:B_\beta(v)=0,\quad
      \iota_v\beta(y)\in\CC\,\ell_v(y)\ell_v\}
\]
を定める。第二条件は、$y$ に関して線形な1形式の空間の中で表示した直線に入る、
という条件である。$v$ の非零スカラー倍によって条件は変わらない。

$v=\partial_y$ の場合、これらの条件と閉性・primitive 性は
\begin{align*}
 \beta={}&t x\dd x\wedge\dd y
 +(b x+c z+d w)\dd x\wedge\dd z\\
 &+(e x+f z+g w)\dd x\wedge\dd w-tx\dd z\wedge\dd w,
 \qquad d-f=t
\end{align*}
を与える。従って $S_{[\partial_y]}=S$、次元は6である。
$\Sp_4$ は $\PP(U)=\PP^3$ に推移的であり、定義は同変なので
\[
 \mathcal I=\{([v],\beta):\beta\in S_{[v]}\}
                  \subset\PP^3\times W
\]
は階数6のベクトル束の全空間である。
条件は $\beta$ に線形で、$\ell_v\ne0$、その相対階数は一定である。
第二条件を二つのテンソルの比例性の小行列式で書けば、$\mathcal I$ はこの積内の閉部分集合となる。

\begin{theorem}[二成分の共通一次像]\label{thm:ambient-image}
$Z=\operatorname{pr}_2(\mathcal I)$ は $W$ の閉既約9次元部分集合であり、
$\rank B_\beta\le3$ に含まれる。
$\mathcal C_A,\mathcal C_B$ に対応する primitive gauge の $S_{61}$ 成分の
一次ジェット像は、複素点の集合としてともに正確に $Z$ である。
\end{theorem}
\begin{proof}
$\mathcal I$ は既約9次元であり、射影 $\PP^3\times W\to W$ は proper なので、
$Z$ は閉既約、$\dim Z\le9$ である。
$B_\beta(v)=0$ より係数階数は3以下。
次の明示的な点を考える。
\[
 \beta_*=x\dd x\wedge\dd y+(z+w)\dd x\wedge\dd z
            +w\dd x\wedge\dd w-x\dd z\wedge\dd w.
\]
これは $S$ に属し、係数行列の行 $(12),(13),(14)$、列 $x,z,w$ の小行列は
\[
 \begin{pmatrix}1&0&0\\0&1&1\\0&0&1\end{pmatrix}
\]
である。従って係数階数は3、核は $\CC\partial_y$。
階数3の非空開部分では $[v]=\PP\ker B_\beta$ が一意かつ代数的に復元されるので、
incidence 射影はその部分で同型となり、$\dim Z=9$ を得る。
独立な確認として、六 parameter と $\Sp_4$ 作用を合わせた微分にも
9次非零小行列式 $-2016$ がある\cite{ambient-certificate}。

推移性から $Z=\Sp_4\cdot S$。
$S$ の全点を両原族で実現したので、$Z$ の全点は両成分から実現する。
逆の包含は閉包の点にも確認する必要がある。
非退化評価開部分 $\mathcal U$ 上で、交代行列を $J$ にする線形 frame は
非零 minor を局所化した Zariski 開集合ごとに正則に選べる。
続く primitive 正規化も多項式である。
$Z$ は $\CSp_4$ 不変かつ閉なので、正規化された一次項が $Z$ に属するという条件は
frame の選択によらず、$\mathcal U$ 上の閉条件となる。
上の任意の基点の計算により、原族の $\PGL_5$ 飽和の $\mathcal U$ との交わりはこの閉条件を満たす。
その閉包である $\mathcal C_A\cap\mathcal U$、$\mathcal C_B\cap\mathcal U$ でも条件は保たれる。
定数項を $J$ に固定して perspective を除いた slice に制限すれば、逆包含を得る。
これは複素点の像の主張であり、scheme-theoretic image や定義イデアルの根基性を追加しない。
\end{proof}

\begin{corollary}
既存の成分次元28と $S_{61}$ の次元差9を用いると、
両者に対応する19次元成分の $Z$ 上の一般 fibre の次元は10である。
\end{corollary}
\begin{proof}
どちらも既約19次元から既約9次元への優越写像なので、一般 fibre の次元定理を用いる。
\end{proof}

この結果は二成分の交叉次元、fibre の既約性・同型類、同じ一次項を持つ構造の同一性を決めない。
二成分の相異性は既存の退化因子の重複度によって保たれ、その違いは一次項より高い係数に残る。

\section{一次資料照合と一般一次項の五つの延長}\label{sec:five-lifts}
\subsection{MPS40の個別図式の照合}
MPS\cite{mps}の定理1.4、Table 2、例6.24は、SNC locus と交わる40成分と
その装飾 smoothing diagram の対応を与える。
内部ID T01--T40を一次資料の全40図式に照合し、辺、対頂点の次数、零 biresidue の辺を
$S_5$ の全120置換で比較した\cite{mps-crosswalk}。
全40個の同型類は相異なり、原本カタログの集合と一致する。
この照合は既刊定理全体の独立再証明ではない。

次数の規約は、行和零の反対称行列 $B$ に対して
\[
 \theta_{ijk}=\frac{b_{jk}+b_{ki}}{b_{ij}}
\]
であり、各辺の三つの次数の和は2となる。
これは Matviichuk\cite{matviichuk}の第2節と一致する。
MPS v2の式(5.11)の印刷された分子は逆符号で、行和零と併用すると和が $-2$ になる。
今回の照合ではこの不一致を明記し、2023年資料の規約と実際の40図式に直接比較した。

period stratum の次元を $s$、smoothable 辺数を $e$ とすると、
MPSの非障害局所模型とトーリック点の $\PGL_5$ 軌道次元20から、
射影成分次元は $19+s+e$ となる。各 period 条件の核を $\QQ$ 上で求め、
原本の基底、CSV、一次図式に表示された次元と照合した。
次元23が29個、24が8個、25が3個である。
25次元の三つは $T01,T39,T40$、Table 2のHTMLでそれぞれ
\texttt{pic40}, \texttt{pic1}, \texttt{pic2}に当たる。

楕円型の名称も\cite{matviichuk}の命題4.2と定理4.3で確認した。
T39の biresidue 例は頂点置換により $\mathcal C_{5,2}$、T40は $\mathcal C_{5,1}$ に等しい。
$k=((1+\widetilde k)^{-1}-1)\bmod5$ により、それぞれ $q_{5,1}$、$q_{5,2}$ を得る。
個別40件の台帳を得ても、それぞれの一般族の S61 座標や定義イデアルの構築は別に必要である。

\begin{corollary}[既知の一般階数四成分の一次像]\label{cor:known-dominant}
既存の51個の一般階数四成分のうち、$W$ へ優越するものはちょうど
$T01,T39,T40$ の三つである。
\end{corollary}
\begin{proof}
この三つの優越性は定理\ref{thm:toric-dominance}、\ref{thm:elliptic-dominance}による。
残るMPS37個は射影次元24以下であり、S61成分次元は15以下だから $W$ に優越しない。
追加11成分のうち $\mathcal C_A,\mathcal C_B$ 以外の9個も射影次元24以下である。
残る二つの19次元 slice 成分の像は、定理\ref{thm:ambient-image}の9次元集合 $Z$ である。
\end{proof}

\begin{theorem}[一般一次項の少なくとも五つの延長]\label{thm:five-lifts}
$W$ のある非空 Zariski 開部分 $W^\circ$ の任意の点について、
全 $\widetilde S_{61}$ の fibre は少なくとも五つの相異なる、孤立した被約点を持つ。
三点は $T01$、少なくとも一点ずつは $T39,T40$ に対応する成分から取れる。
\end{theorem}
\begin{proof}
対応する被約成分を $C_1,C_2,C_3$ とする。
各 $C_i\to W$ は16次元既約空間の間の優越写像だから一般に有限であり、
像は非空開部分を含む。
定理\ref{thm:etale}により、全スキームが滑らかで射影が \'etale となる
非空開部分 $C_i^{\rm et}\subset C_i$ がある。
補集合の次元は高々15であるから、その像の閉包を $W$ から除くことができる。
三つの優越像に含まれる開部分を取り、これらの閉集合をすべて除く。

さらに $i\ne j$ なら相異なる16次元既約成分の交叉 $C_i\cap C_j$ は次元15以下である。
その一次部分像の閉包もすべて除けば、なお非空開部分 $W^\circ$ が残る。
さらに定理\ref{thm:toric-degree}の有限\'etaleな開集合へ縮めると、
各 $\beta\in W^\circ$ にはT01から三つの相異なる逆像がある。
T39とT40からもそれぞれ少なくとも一つの逆像があり、
交叉を避けるので合わせて少なくとも五点は相異なる。
全スキームの射影がその点で \'etale であるから、各点は fibre の孤立した被約点となる。
\end{proof}

五点は、同じ正規化座標と一次項を持つ相異なる射影 Poisson 双ベクトルを意味する。
$\PGL_5$ 同型類の個数を数えたものではない。
T01の一般次数は3であるが、T39・T40の一般次数は未決であり、未知成分が $W$ を支配する可能性も残る。
特に全 fibre が有限、またはちょうど五点であるという結論は得られない。

\section{未飽和斉次化の無限遠に現れる旗}\label{sec:infinity}
固定一次部分 $\beta\in W$ の全S61ファイバーは45個の重み2座標qに関する134二次式
$F_i(q)=0$ で与えられる。各式を個別に二次斉次化して
\[
 X_\beta^h=V(F_1^h,\ldots,F_{134}^h)\subset\PP^{45}_{[q:t]}
\]
と置く。これは、斉次化変数tで飽和した真の射影閉包とは区別する。
この節は\cite{infinity-certificate}の障害を示し、アフィンファイバーの正次元性は主張しない。

\begin{theorem}[旗による無限遠の閉部分]\label{thm:infinity}
$\Omega(v,w)=v^tJw$ とする。独立なv,wで $\Omega(v,w)=0$ を満たすものに対し
\[
 Q_{v,w}(y)=\ell(y)^2v\wedge w,\qquad\ell=-Jv
\]
と置く。isotropic flag $\langle v\rangle\subset\langle v,w\rangle$ から射影方向
$[Q_{v,w}]$ への写像は閉埋め込みであり、
\begin{equation}\label{eq:infinity-incidence}
 [Q_{v,w}]\in X_\beta^h\cap\{t=0\}
 \quad\Longleftrightarrow\quad T_\beta(v,w):=\beta_v(v,w)=0.
\end{equation}
βが非零ならこの旗の部分は純3次元であり、βが零なら4次元旗多様体全体である。
\end{theorem}
\begin{proof}
$\ell(v)=\ell(w)=0$ と $J\ell=v$ から $[J,Q]=[Q,Q]=0$。
従ってQは45次元二次方向空間の非零点を与える。
旗の写像は二次Veronese写像と平面のPl\"ucker埋め込みのSegre積の制限であり、閉埋め込みである。
Qの共通二次因子 $\ell^2$ と定数双ベクトルから、線と平面も直接回復できる。
旗多様体 $\mathcal F=F_{\rm iso}(1,2;4)$ は滑らかな既約4次元多様体である。

代表 $v=-e_2,w=-e_4$ では $Q=y_1^2\partial_2\wedge\partial_4$。
旧S61座標では $u_{38}=1$、他の重み2座標は零となる。
16個の重み1座標を自由のまま、134式のqに関する最高二次部分へ代入すると、
非零なのは0始まりの第84式だけで、その値は $2u_4$ となる。
同じ座標で $T_\beta(-e_2,-e_4)=-u_4$ だから、代表では\eqref{eq:infinity-incidence}を得る。

この代表の計算を群作用で移す根拠を明示する。
固定LでS61の低次Jacobi条件を解く再構成は、45個の二次座標に関してアフィン線形である。
実際、保存された55消去式の重み2変数に関する次数は高々1である。
有限perspective公式\eqref{eq:Qtransform}はQをLに依存する項だけ平行移動するので、
二次方向そのものを変えず、三次方向を\eqref{eq:Ctransform}に従って移す。
原点固定の局所同型はJacobiの原点での消滅次数を保つため、低次Jacobiを解いた空間は
同じ再構成空間へ移る。再構成を $P_0+sB$ と書けば最高二次条件は $[B,B]=0$ であり、
pushforwardで保たれる。154式から134式へのBianchi除去の20恒等式は、係数に定数または
重み1変数しか使わないため、固定Lで最高二次部分を取っても同値性を保つ。
以上の性質は保存証明書でも検算した。

さらにperspective変換による一次項の変化は
\[
 T_{\beta'}(v,w)-T_\beta(v,w)=-3a(v)\Omega(v,w)
\]
であり、isotropic pair上では零となる。
従って古いgaugeからprimitive gaugeへ移しても条件は同じである。
$\Sp_4$ は $\mathcal F$ に推移的に作用するので、代表の必要十分条件は全旗へ移る。

線 $\CC v$ を固定すると、wの選択は $\PP(v^\perp/\CC v)\simeq\PP^1$ であり、
Tの条件は二次元ベクトル空間上の一つの線形条件である。従って各線上に必ず解がある。
非零βについて、Tが旗全体で零にならないことを示す。
$\alpha=\iota_E\beta$ と置くと、係数次数1と閉性からCartan公式は $\dd\alpha=3\beta$ を与える。
Tがすべてのisotropic pairで零なら、既約な双線形式 $\Omega(v,w)$ による可除性と
双次数 $(2,1)$ から、ある一次形式fについて
\[
 \alpha=f\iota_E\Omega,\qquad
 \beta=\tfrac13(\dd f\wedge\iota_E\Omega+2f\Omega)
\]
となる。primitive contractionを取ると右辺は $(f+4f)/3=5f/3$ なのでf=0、従ってβ=0。
これは仮定に反する。よってβが非零ならTの零点は $\mathcal F$ 上の非零な有効Cartier因子であり、
純3次元となる。β=0なら全旗が入る。
\end{proof}

切断の単射性は、primitive閉一次2形式の16基底から作るTと
$\Omega(v,w)v_i$ の4多項式の計20個の係数行列でも点検した。
20次小行列式は625である。別の有理旗でも全16一次変数を残し、
134個の最高二次残余の線形包がTの張る一次元空間と一致することを確認した。
これらの追加検算を、上のgauge・群作用の証明の代わりには用いない。

\begin{remark}[一般有限性への障害の範囲]
定理\ref{thm:infinity}により、未飽和の $X_\beta^h$ はどの一次部分でも零次元ではない。
従って、この模型の零次元特殊化を用いて全S61ファイバーの一般有限性を示す方法は使えない。
真の射影閉包では $\bigl(F_i^h\bigr):t^\infty$ を取る必要があり、
上記の旗の部分はその操作で消え得る。
この結果はアフィンファイバーの正次元部分の存在や、未知成分の存在を証明しない。
次の手段は、飽和を制御するsyzygy、または明示した開条件でのアフィン消去である。
\end{remark}

\section{検証器と再現範囲}
\begin{longtable}{@{}p{58mm}p{98mm}@{}}
\toprule
検証器 & 検証する内容\\\midrule\endhead
\nolinkurl{code/s61_weight1_audit.py} &
旧実装を import しない17項目の有理数体上の確認。
perspective の有限行列恒等式、radial 三次項の相殺、trace 変換、
閉性と primitive 条件の階数、冪等射影を扱う。\\
\nolinkurl{code/weight1_structure.py} &
閉 primitive 一次2形式の16次元核、Lie 作用の不変性、
上記の一般軌道の小行列式、階数一代表の Jacobi と Pfaffian を扱う。\\
\nolinkurl{code/weight1_rank_one_fiber.py} &
保存された S61 の全134式を固定一次項へ特殊化し、
定数係数線形包と固定 fibre の接空間を exact に計算する。\\
\nolinkurl{code/toric_weight1_dominance.py} &
六変数 toric 族の Jacobi、線形 frame、primitive 一次項を恒等的に確認し、
16次小行列式から優越性を証明する。\\
\nolinkurl{code/verify_s61_families.py} &
一般 toric 係数族と almost-abelian 族の S61 座標写像を有理関数恒等式として確認する。
これらの族の像を成分全体の定義イデアルと同一視しない。\\
\nolinkurl{code/elliptic_weight1_dominance.py} &
二つの楕円族の一般 Jacobi、符号比較、frame 微分、一次射影の16次小行列式。\\
\nolinkurl{code/s61_known_fiber_audit.py} &
三点の全134式、復元、独立 Jacobi、全61列と重み2の45列の Jacobian。\\
\nolinkurl{code/toric_coframe_fiber.py} &
coframe 復元の一般 Jacobi・Pfaffian、一次式、16次小行列式、再ラベルの点検。\\
\nolinkurl{code/toric_algebra_fiber.py} &
20変数・96二次式、一次写像の右逆と核、判別式、全双ベクトルとノルムの復元。\\
\nolinkurl{code/toric_algebra_fiber_review.py} &
一般40係数からの独立導出、結合性の省略ブロックの生成イデアル所属、文字からの復元。\\
\nolinkurl{code/toric_algebra_degree.py} &
普遍的18一次帰結、有限近傍のleading行列、全特殊ファイバーの三点、判別式とB\'ezout証明書。\\
\nolinkurl{code/toric_degree_review.py} &
二次係数行列の左核からの独立消去、全残余イデアルと三次式、半単純性、元証明書との照合。\\
\nolinkurl{code/s61_infinity_flags.py} &
代表旗の全一次変数を残す最高二次残余、gaugeとBianchiの次数、切断の単射性、別の旗の検算。\\
\nolinkurl{code/mps_catalog_audit.py} &
一次資料の40図式と内部IDの全単射、零辺・period 条件・次元・楕円型名称の照合。\\
\bottomrule
\end{longtable}

独立検証は $\QQ$ 上、SymPy 1.14.0を用いる。
有限体の標本や浮動小数の階数を、標数零の既約分解の代わりには使っていない。
再現時は保存済み証明書を上書きせず、新しい出力名を指定する。
\begin{verbatim}
.venv/bin/python code/s61_weight1_audit.py \
  --output reports/new-run/perspective.json
.venv/bin/python code/weight1_structure.py \
  --output reports/new-run/weight1.json
.venv/bin/python code/weight1_rank_one_fiber.py \
  --output reports/new-run/rank-one-fiber.json
\end{verbatim}

\section{残る問題}
一般一次項の延長可能性、T01の一般次数3、既知三成分からの少なくとも五つの相異なる延長は証明した。
次の課題はT01の分岐・例外層、三次被覆のモノドロミー、T39・T40の一般次数、
および未知成分からの延長を排除できるかである。
楕円型成分については一般 fibre の明示的な表示も未決である。
全S61の未飽和斉次化には定理\ref{thm:infinity}の旗の部分が必ず現れるため、
斉次化変数による飽和の制御、またはアフィンでの局所化と消去が必要となる。

特殊層では、係数階数一の固定 fibre $F$ の極小素イデアル分解、成分次元、
既知の全 Poisson 成分との対応が未決である。
階数7の minor が非零の近傍での局所消去や、連結安定化群 $H$ の作用が次の削減候補となる。
$\mathcal C_A,\mathcal C_B$ は共通の閉集合 $Z$ 上に10次元の一般 fibre を持つが、
その既約性・同型類や二成分の交叉は決定していない。
係数階数2,3,4の全層と、toric 境界の全解の分類も残る。

全分類の第一段階には、$S_{61}$ のすべての極小成分を尽くす証明が必要である。
本稿はその上界を与えず、下界57を正確な総数へ変更しない。
特殊軌道・成分交叉・冪零構造は、さらに別の未解決目標として維持する。

\begingroup\small
\begin{thebibliography}{99}
\bibitem{mps} M. Matviichuk, B. Pym, T. Schedler,
\textit{A local Torelli theorem for log symplectic manifolds},
arXiv:2010.08692v2 (2021), Theorem 1.4, Table 2, Example 6.24.
\url{https://arxiv.org/abs/2010.08692v2}.
\bibitem{matviichuk} M. Matviichuk,
\textit{Elliptic log symplectic brackets on projective bundles},
arXiv:2310.05284v1 (2023), Section 2, Proposition 4.2, Theorem 4.3.
\url{https://arxiv.org/abs/2310.05284v1}.
\bibitem{ortp} G. Ortenzi, V. Rubtsov, S. R. Tagne Pelap,
\textit{On the Heisenberg invariance and the Elliptic Poisson tensors},
arXiv:1001.4422v2 (2010), equations (16), (18), (19), (21).
\url{https://arxiv.org/abs/1001.4422v2}.
\bibitem{v2} Poisson P4 研究資料、第2版、
\nolinkurl{legacy/v02/p4_poisson_completeness_v2.tex}, 第3節。
61変数模型、三角消去、Bianchi 恒等式の原資料。
\bibitem{perspective-note} 同プロジェクト、
\nolinkurl{research/notes/S61_MODEL_AUDIT_20261030.md}、
および\nolinkurl{reports/20261030-s61-weight1/finite_identity_certificate.json}。
\bibitem{weight1-certificate} 同プロジェクト、
\nolinkurl{research/notes/S61_WEIGHT1_20261030.md}、
および\nolinkurl{reports/weight1_structure_20261030.json}。
\bibitem{fiber-certificate} 同プロジェクト、
\nolinkurl{reports/weight1_rank_one_fiber_20261030_final.json}。
入力ファイルの SHA-256、45変数の式、行列と小行列式を含む。
\bibitem{toric-certificate} 同プロジェクト、
\nolinkurl{research/notes/TORIC_WEIGHT1_DOMINANCE_20261030.md}、
\nolinkurl{reports/TORIC_WEIGHT1_DOMINANCE_20261030.json}、
\nolinkurl{code/toric_weight1_dominance.py}。
\bibitem{ambient-source} 同プロジェクト、第6版、
\nolinkurl{legacy/v06/ambient_families.tex}、二つの28次元成分の構成と証明。
\bibitem{ambient-certificate} 同プロジェクト、
\nolinkurl{research/notes/AMBIENT_WEIGHT1_IMAGE_20261030.md}、
\nolinkurl{reports/AMBIENT_WEIGHT1_IMAGE_20261030.json}、
\nolinkurl{code/ambient_weight1_image.py}。
\bibitem{component-catalog} 同プロジェクト、
\nolinkurl{legacy/v02/baseline/data/p4_toric_components_40.csv}、
\nolinkurl{research/COMPONENTS.json}。既存成分の同定・次元の台帳。
\bibitem{elliptic-certificate} Poisson P4 研究プロジェクト、
\nolinkurl{research/notes/ELLIPTIC_WEIGHT1_20261031.md}、
\nolinkurl{reports/ELLIPTIC_WEIGHT1_DOMINANCE_20261031.json}、
\nolinkurl{code/elliptic_weight1_dominance.py}。
\bibitem{known-fiber-certificate} 同プロジェクト、
\nolinkurl{research/notes/S61_KNOWN_FIBER_AUDIT_20261031.md}、
\nolinkurl{reports/S61_KNOWN_FIBER_AUDIT_20261031.json}、
\nolinkurl{code/s61_known_fiber_audit.py}。
\bibitem{coframe-certificate} 同プロジェクト、
\nolinkurl{research/notes/TORIC_COFRAME_FIBER_20261031.md}、
\nolinkurl{reports/TORIC_COFRAME_FIBER_20261031.json}、
\nolinkurl{code/toric_coframe_fiber.py}。
\bibitem{algebra-certificate} 同プロジェクト、
\nolinkurl{research/notes/TORIC_ALGEBRA_FIBER_20261031.md}、
\nolinkurl{reports/TORIC_ALGEBRA_FIBER_20261031.json}、
\nolinkurl{code/toric_algebra_fiber.py}。
\bibitem{algebra-review} 同プロジェクト、
\nolinkurl{research/notes/TORIC_ALGEBRA_FIBER_REVIEW_20261031.md}、
\nolinkurl{reports/TORIC_ALGEBRA_FIBER_REVIEW_20261031.json}、
\nolinkurl{code/toric_algebra_fiber_review.py}。
\bibitem{infinity-certificate} 同プロジェクト、
\nolinkurl{research/notes/S61_INFINITY_FLAGS_20261031.md}、
\nolinkurl{reports/S61_INFINITY_FLAGS_20261031_final.json}、
\nolinkurl{code/s61_infinity_flags.py}。
\bibitem{degree-certificate} 同プロジェクト、
\nolinkurl{research/notes/TORIC_GENERIC_DEGREE_20261031.md}、
\nolinkurl{reports/TORIC_ALGEBRA_DEGREE_20261031.json}、
\nolinkurl{code/toric_algebra_degree.py}。
\bibitem{degree-review} 同プロジェクト、
\nolinkurl{research/notes/TORIC_GENERIC_DEGREE_REVIEW_20261031.md}、
\nolinkurl{reports/TORIC_DEGREE_REVIEW_20261031.json}、
\nolinkurl{code/toric_degree_review.py}。
\bibitem{mps-crosswalk} 同プロジェクト、
\nolinkurl{research/notes/MPS40_CROSSWALK_20261031.md}、
\nolinkurl{reports/MPS40_CROSSWALK_20261031.json}、
\nolinkurl{code/mps_catalog_audit.py}。
\end{thebibliography}
\endgroup
\end{document}
